\chapter[Recursive Relational Geometry]{Recursive Relational Geometry: An Expository Track}

\begin{remark}[Epistemic status]
This appendix is intentionally separated from the theorem-bearing $\ell=5$ audit. It contains one established external mathematical result, several framework definitions, and a set of interpretive propositions about scale and relational description. The definitions are modeling choices. The interpretations are not empirical or ontological conclusions. In particular, recursive geometric constructions do not establish that physical reality is recursively packed, that nature is fundamentally fractal, or that consciousness follows from nested organization.
\end{remark}

\section{Relational Worlds Across Scale}

Reality can be described at many scales without requiring those scales to be separate realities. A system that appears to be one organized whole at one resolution may reveal itself as a network of interacting parts at another; that network may itself participate as one component inside a larger organization. The purpose of a recursive relational model is not to declare which scale is fundamental, but to make explicit what changes when the descriptive scale changes and what relational organization, if any, persists.

A useful hierarchy is

\[
\begin{gathered}
\text{vertices}
\longrightarrow
\text{graphs}
\longrightarrow
\text{domains}\\
\longrightarrow
\text{relations among domains}
\longrightarrow
\text{recursive organization}.
\end{gathered}
\]

These arrows indicate levels of representation, not a causal or metaphysical production chain. A structure can be a whole relative to the relations it contains and a part relative to the relations in which it participates.

\section{Framework Definition: A Relational Domain}

\begin{definition}[Relational domain]
A relational domain is a modeling tuple
\[
C_i=(c_i,r_i,G_i),
\]
where $c_i\in\mathbb R^2$ is a chosen spatial center, $r_i>0$ is a characteristic geometric scale, and
\[
G_i=(V_i,E_i)
\]
is an internal graph describing relations among the components assigned to that domain.
\end{definition}

The geometric variables $(c_i,r_i)$ and the internal graph $G_i$ carry different information. The first describe where and at what scale the domain is represented; the second describes internal organization. The graph can be encoded by an adjacency matrix $A_i$ with

\[
(A_i)_{ab}=
\begin{cases}
1,&\text{if vertices }a\text{ and }b\text{ are related},\\
0,&\text{otherwise}.
\end{cases}
\]

Simple graph statistics include

\[
N_i=|V_i|,
\qquad
M_i=|E_i|,
\qquad
\bar d_i=\frac{2M_i}{N_i}.
\]

These quantities measure aspects of internal graph structure. They do not by themselves determine what the edges mean physically, biologically, informationally, or socially.

\section{Relations Within and Relations Between}

A collection of relational domains can itself form a higher-order graph

\[
T=(\{C_i\},E_T).
\]

An edge $(C_i,C_j)\in E_T$ may encode adjacency, contact, coupling, information exchange, causal accessibility, functional dependence, or some other explicitly declared relation between complete domains. The total relational structure can then be written schematically as

\[
E_{\mathrm{total}}
=
\left(\bigcup_iE_i^{\mathrm{int}}\right)
\cup E^{\mathrm{ext}}.
\]

The common graph grammar does not imply that an edge inside a domain and an edge between domains represent the same physical interaction. A molecular relation, a cellular relation, an organism-level relation, and an institutional relation can all be represented graphically without becoming semantically identical. Formal similarity is not ontological identity.

\section{Established Example: Recursive Circle Packing}

A particularly clean mathematical example of recursive refinement is supplied by Apollonian circle packing. Repeatedly filling curvilinear gaps among mutually tangent circles generates circles at progressively smaller scales \cite{graham2003Apollonian}.

If the curvature of a circle is

\[
k_i=\frac1{r_i},
\]

then four mutually tangent circles satisfy Descartes' relation.

\begin{theorem}[Descartes' Circle Theorem]
For four mutually tangent circles with signed curvatures $k_1,k_2,k_3,k_4$,
\[
(k_1+k_2+k_3+k_4)^2
=
2(k_1^2+k_2^2+k_3^2+k_4^2).
\]
\end{theorem}

Given three mutually tangent circles, the relation constrains the curvature of a fourth circle tangent to them. Repetition creates additional gaps and additional tangent circles. Abstractly one may write a refinement operation as

\[
\mathcal P_{n+1}=\mathcal R(\mathcal P_n),
\]

with characteristic scales satisfying an idealized hierarchy

\[
r_0>r_1>r_2>\cdots>0.
\]

The established mathematical conclusion is narrow and sufficient:

\begin{quote}
A finite local rule can generate recursively finer structure without requiring a different rule at every scale.
\end{quote}

That existence result does not license the stronger statement that the physical universe has an Apollonian ontology. The packing is an expository model of recursive refinement, not evidence that reality literally consists of nested tangent domains.

\section{Framework Definition: Recursive Relational State}

\begin{definition}[Recursive relational state]
A recursive relational model may be represented as
\[
\mathcal X=
\left(
\{C_i=(c_i,r_i,G_i)\},
T,
\mathcal R,
\Phi
\right),
\]
where $G_i$ describes internal relations, $T$ describes relations among domains, $\mathcal R$ specifies a refinement rule, and $\Phi$ is a chosen map extracting the organizational form whose persistence is being studied.
\end{definition}

Nothing in this definition asserts that $\mathcal X$ is an ontologically fundamental state. It is a bookkeeping structure useful when a scientific problem contains nested organizations and explicit scale changes.

\section{Resolution Does Not Create Difference}

Let

\[
\Lambda_n:\mathscr R\longrightarrow\mathcal L_n
\]

represent a descriptive language at resolution $n$. It induces the equivalence relation

\[
x\sim_n y
\quad\Longleftrightarrow\quad
\Lambda_n(x)=\Lambda_n(y).
\]

If a richer descriptive language preserves distinctions absent from the coarser one, then

\[
[x]_{n+1}\subseteq[x]_n.
\]

Strict inclusion,

\[
[x]_{n+1}\subsetneq[x]_n,
\]

means that states identified by the coarser language are separated by the richer one. The richer description does not manufacture the difference; it acquires enough resolution to represent it.

This is an expository reuse of the descriptive-refinement principle already proved and illustrated in the $\ell=5$ invariant audit:

\[
\begin{gathered}
\text{restricted descriptive indistinguishability}\\
\not\Rightarrow\\
\text{identity under every richer description}.
\end{gathered}
\]

The recursive model adds no new theorem to that statement. It supplies a scale-based vocabulary for applying the same distinction to nested systems.

\section{Coherence Across Scale}

Let $x_n$ denote a realization at scale $n$ and let $\Phi(x_n)$ denote the organizational form selected by the model. Exact cross-scale persistence would satisfy

\[
x_n\neq x_{n+1}
\]

while

\[
\Phi(x_n)=\Phi(x_{n+1}).
\]

A practical formulation allows

\[
d_{\mathcal F}(\Phi(x_n),\Phi(x_{n+1}))\le\epsilon.
\]

This is not a new theorem. It is a scale-indexed application of the book's standing definition of relational coherence:

\begin{quote}
Coherence is the persistence of defining relationship through variation.
\end{quote}

The important methodological point is that self-similarity is only one possible realization of cross-scale coherence. Two scales need not look visually identical. They need only preserve the declared relation or organizational invariant relevant to the scientific question.

\section{Interpretive Proposition: Whole and Part Are Scale-Relative Roles}

A domain $C_i$ can be treated as a whole relative to its internal graph $G_i$ and simultaneously as a node in the higher-order graph $T$. This motivates the interpretive proposition

\begin{quote}
Whole and part are often scale-relative relational roles rather than absolute labels attached permanently to an object.
\end{quote}

The proposition is useful in multiscale science but should not be mistaken for a theorem of ontology. A cell can be modeled as an organized whole relative to molecular interactions and as a component relative to tissue organization. A tissue can be modeled as a whole relative to cells and as a component relative to an organ. The mathematical representation makes the change of scale explicit; it does not prove that every natural hierarchy has the same architecture.

\section{Interpretive Proposition: Finer Scale Is Not Greater Ontological Depth}

Recursive refinement can tempt the reader to equate smaller scale with deeper reality. Nothing in the mathematics warrants that inference. A microscopic description may contain more variables while being less explanatory for a macroscopic question. A coarse description may discard vast amounts of microscopic information while preserving the organization relevant to the phenomenon under study.

The appropriate methodological statement is therefore

\begin{quote}
Finer descriptive resolution does not, by itself, imply greater ontological depth.
\end{quote}

The scientifically useful question is not simply how far a system can be subdivided, but at what resolution the relations required by the phenomenon become expressible.

\section{Internal and External Coherence}

The expository model also permits two distinct coherence functions. Let

\[
C_{\mathrm{int}}(C_i)
\]

measure organization internal to a domain and let

\[
C_{\mathrm{ext}}(C_i,T)
\]

measure the stability of its relations to the larger graph. A system can be internally coherent while weakly integrated with its surroundings, or can undergo substantial internal turnover while preserving a stable external role. The distinction illustrates why the ordinary word \emph{identity} can conceal multiple equivalence relations: internal form, causal history, external function, dynamical role, and symmetry class need not coincide.

\section{Resolution Is Not Orientation}

Recursive enrichment can reveal additional components and relations without determining their direction. A graph may contain a relation between $A$ and $B$ without determining whether the relevant dependence is

\[
A\longrightarrow B,
\qquad
A\longleftarrow B,
\]

or

\[
A\longleftarrow R\longrightarrow B.
\]

Thus a richer recursive description may answer what additional relations become visible while leaving open which way those relations run. The distinction remains

\[
\text{resolution}\neq\text{orientation}.
\]

And the cross-cutting rule of the main text still applies:

\[
\text{relation}\not\Rightarrow\text{direction of predication}.
\]

Recursive geometry therefore enriches the expository language of form and resolution. It does not settle the later causal or ontological audit.

\section{What the Model Is Good For}

The recursive-relational framework is useful when a problem genuinely contains nested organizations, changing descriptive scales, or systems that are simultaneously wholes and components. It can help a modeler ask four disciplined questions:

\begin{enumerate}[label=\arabic*.]
\item What becomes distinguishable when the descriptive scale is refined?
\item Which organizational relations persist through that refinement?
\item Which relations are internal to the selected system boundary, and which connect the system to a larger organization?
\item Which apparent arrows are properties of the evidence, and which were introduced by the descriptive grammar?
\end{enumerate}

Its proper conclusion is therefore methodological rather than ontological. A world may be modeled as containing networks within networks, and a finite recursive rule may generate indefinitely fine mathematical structure. What follows is a vocabulary for multiscale relational analysis. What does not follow is a proof that reality itself is fundamentally recursive, fractal, Apollonian, conscious, or governed by any particular metaphysical hierarchy.
